\documentclass{report}
\usepackage{amsmath,amssymb}
\setlength{\parindent}{0mm}
\setlength{\parskip}{1em}
\begin{document}
\begin{center}
\rule{6in}{1pt} \
{\large Andrey Prokopenko \\
{\bf Energy minimization in parallel setting}}

Sandia National Laboratories \\ PO Box 5800 MS 1320 \\ Albuquerque \\ NM 87185
\\
{\tt aprokop@sandia.gov}\\
Jeremie Gaidamour\\
Jonathan Hu\\
Ray Tuminaro\end{center}

The linear systems arising from multi-physics simulations may present
significant challenges for traditional algebraic multigrid (AMG) solvers.
AMG methods based on energy minimization prove to be an effective
approach to some of these problems. The flexibility of these methods is
based on two components: an {\it a priori} chosen nonzero pattern of a
prolongator, and interpolation of low energy modes, both of which act as
constraints in an optimization problem. The parallel solution of this
constrained problem presents some practical challenges, e.g. how to
handle situations where the problem is over-constrained, or when the
nonzero pattern is non-local. In this talk, we discuss a parallel
implementation of energy minimization AMG, developed using a flexible
parallel multigrid framework called MueLu, part of the Sandia Trilinos
project. We discuss the challenges, and demonstrate that the resulting
algorithms are effective and scalable on challenging problems.


\end{document}
